% !TeX root = ./0-PhD-Thesis-JLBG.tex
\label{cap: introduccion}
\section{Background}
Since the mid-20th century, various techniques have been proposed for manipulating objects at the micro and nanoscale \cite{ashkin1970acceleration, ashkin1987optical, Ashkin, custance2009atomic, dholakia2011shaping, marago2013optical, romo2020controlled}. Optical tweezers, based on the electromagnetic forces produced by focused light beams, have been widely employed for trapping and manipulating microobjects, including viruses and bacteria, thereby significantly impacting technological development, particularly in the field of medicine \cite{ashkin1970acceleration, ashkin1987optical, Ashkin}.

In 2004, García de Abajo introduced the possibility of manipulating nanoobjects using transmission electron microscopes (TEMs) \cite{GarciadeAbajo0}. Subsequently, experimental evidence has shown that TEMs induce motion and rotation in nanoparticles (NPs) \cite{Batson01, zheng2012electron}, leading to the development of a manipulation technique termed \textit{electron tweezers} \cite{Batson, oleshko2005electron, Oleshko}, referencing optical tweezers.

In various experimental studies on electron tweezers, it has been observed that the transfer of angular and linear momentum from the electron beam to a nanoparticle (NP) depends on both the speed of the electrons in the beam and the impact parameter \textemdash the effective distance between the trajectory of the electron beam and the center of the NP \cite{OLESHKO2013203, Oleshko, Batson, Batson01, zheng2012electron, xu2010transmission}. By modifying the impact parameter, it is possible to induce either attractive or repulsive interaction between the electron beam and the NP, and it is also feasible to alter the direction of the induced rotation on the NP \cite{OLESHKO2013203, Batson, Oleshko}.

The scanning transmission electron microscope (STEM) forms images by employing focused electron beams that scan the area of interest, interacting with the sample \cite{Batson}. Figure \ref{fig: Batson STEM} illustrates a pair of gold particles, one large and one small. The small NP is scanned by the beam of a STEM within the region outlined in white. The STEM electron beam scans line by line, and during scanning, the electron beam remains stationary 20\% of the time at the beginning of each line, awaiting a synchronization signal, resulting in a net current of electrons traveling away from the NP. This current is depicted as a shaded blue region in Figure \ref{fig: Batson STEM}. Therefore, although the electron beam scans the entire area of interest, it can be considered an effective impact parameter with respect to the center of the NP.

\begin{figure}[h!]
\centering
\includegraphics[width=0.6\linewidth]{17-imagenes/1-Intro/STEM}
\caption{\label{fig: Batson STEM} Schematics of the interaction within a STEM between an electron beam and a pair of gold nanoparticles, reproduced and adapted from Ref. \cite{Batson}.}
\end{figure}

Figure \ref{fig: Batson momentum transfer} presents results reported in Ref. \cite{Batson} showing six STEM images of a gold NP with a diameter of $\sim 1.5$ nm in the presence of another NP with a diameter of $\sim 5$ nm, at different times. In all images of Figure \ref{fig: Batson momentum transfer}, the effective impact parameter is located to the left of the NPs (near the left edge of the image). The top three images of Figure \ref{fig: Batson momentum transfer}, taken with an effective impact parameter of $\sim 4.5$ nm, demonstrate an attractive interaction between the beam and the small NP (the large NP is used as a reference for displacement), as it can be observed that the NP approaches the left edge of the images, crossing the vertical white dashed line placed as a visual aid. Conversely, in the bottom images of Figure \ref{fig: Batson momentum transfer}, the effective impact parameter is $\sim 1$ nm, and it is observed that the small NP moves away from the beam, crossing the white dashed line in the opposite direction, thus approaching the right edge of the image, indicating a repulsive interaction. Moreover, using the guide lines drawn on the facets of the large NP, forming a polygon, it can be appreciated that in the top three images, the large NP rotates clockwise, while in the bottom ones, upon changing the impact parameter, it rotates counterclockwise.

\begin{figure}[h!]
\centering
\includegraphics[width=0.6\linewidth]{17-imagenes/1-Intro/Batson}
\caption{\label{fig: Batson momentum transfer} Results reproduced and adapted from Ref. \cite{Batson} showing linear and angular momentum transfer from an electron beam in a STEM to gold nanoparticles, supported on an amorphous carbon substrate.}
\end{figure}

Although the momentum transfer from the electron beam to NPs is observed in the experiment, the development of electron tweezers technique require a theoretical understanding of the problem. Electron beams in a STEM can reach kinetic energies of $400$ keV, with an electric current on the order of pA, equivalent to a pulse of fast electrons traveling at a constant speed of up to $0.83 \, c$ (where $c$ is the speed of light). From this, it is deduced that the emission time of each electron is $\sim 10^{-8}$ s. Since the lifetime of excitations within a metal is typically $\sim 10^{-14}$ s \cite{quijada2010lifetime}, it can be assumed that the NP interacts with one electron at a time \cite{de1999relativistic, GarciadeAbajo-1, deabajo2021optical}. Electron beams in STEM studies have been observed to deviate from a straight path by angles on the order of milliradians \cite{deabajo2021optical, Rivacoba1, krehl2018spectral}, so they can be considered to move in a straight line as long as the electrons travel outside the NP. Therefore, the trajectory of the fast electron can be modeled as $\vv{r}(t)=(b, 0, vt)$, where $v$ is the electron speed and $b$ is the distance between the center of the NP and the electron trajectory (impact parameter), as shown in Figure \ref{fig: system mine}. The electromagnetic response of the NP is modeled by its dielectric function $\epsilon(\omega)$.

From the perspective of quantum mechanics, an electron is considered to be a point-like particle with a fundamental electric charge. However, due to vacuum fluctuations, it has been estimated to have a finite size on the order of its Compton wavelength, which is approximately $\lambda_C = 2.424 \times 10^{-3}$ nm. In the case of a fast electron, such as those traveling within a STEM, its de Broglie wavelength is on the order of $\lambda_B \sim 10^{-2}$ nm \cite{castellanos2021phdthesis}. Therefore, in this work, distances from the NP to the electron greater than $\lambda_C$ and $\lambda_B$ are considered. In a previous work \cite{deabajo2021optical}, the interaction between electron beams and spherical NPs was addressed from the standpoint of quantum mechanics. Remarkably, it was concluded that for distances from the NP to the electron greater than $\lambda_C$ and $\lambda_B$, a classical description of the problem is sufficient.

\begin{figure}[ht!]
\centering
\includegraphics[width=0.5\linewidth]{17-imagenes/1-Intro/system_NP.pdf}
\caption{Nanoparticle with radius $a$, centered at the origin, modeled using a dielectric function $\epsilon(\omega)$, in vacuum. The trajectory of the electron (red dot), with impact parameter $b$ and traveling at constant speed $v$, is shown as a dashed green line. \label{fig: system mine}}
\end{figure}

The interaction between electron beams and spherical NPs has also been studied from the perspective of classical electrodynamics in previous works \cite{GarciadeAbajo0, PRBCoronado, Lagos2, Batson2, xu2010transmission}. The aforementioned studies have focused on calculating linear momentum transfer by solving Maxwell's equations in the frequency domain. To achieve this, a multipolar expansion has been employed to separate the electric and magnetic contributions to the interaction, as well as the contribution of each multipolar order. However, caution is necessary when choosing the dielectric function to model the electromagnetic response of NPs. Since the dielectric function is experimentally measured over a finite frequency range, it is necessary to extrapolate and interpolate it to perform the integration of the Maxwell stress tensor over the entire frequency space. Without sufficient care, this process may result in a non-causal dielectric function, meaning it does not satisfy the Kramers-Kronig relations. Although previous studies \cite{GarciadeAbajo0, PRBCoronado, Lagos2, Batson2, xu2010transmission} successfully reproduced the repulsive behavior of the interaction, recent studies have shown that these works obtained non-physical results by modeling the electromagnetic response of the NP using non-causal dielectric functions \cite{castrejon2021effects, castrejon2021phdthesis}. In these recent works, it is demonstrated that if the non-causal behavior of the dielectric functions is eliminated, the experimentally reported repulsive interaction does not appear. It is interesting to note that in Ref. \cite{castrejon2021phdthesis}, the integrals in frequency space are solved semi-analytically, which were previously solved numerically in Refs. \cite{GarciadeAbajo0, PRBCoronado, Lagos2, Batson2, xu2010transmission}. This allows for an exact understanding of the contribution in frequency space of each multipole, electric or magnetic, to the linear momentum transfer, thereby enabling the calculation of linear momentum transfer from fast electrons to large NPs (up to $a=50$ nm).

The electron tweezers technique will benefit from a detailed theoretical study of angular momentum transfer (AMT). Previous works have discussed angular dynamics briefly (see, for example, Ref. \cite{GarciadeAbajo-1}), but recent studies have enabled its calculation in small NPs, up to $a=5$ nm, using two different methods. The first method models the electromagnetic response of the NP as a point dipole $\vv{p}$, using the polarizability tensor, which is valid only for small NPs \cite{castellanos2021angular}. The second method numerically solves the surface integrals of the Maxwell stress tensor \cite{castellanos2023theory, castellanos2021phdthesis} in frequency space. However, numerical computation limits the size of NPs that can be studied (smaller than 5 nm, for realistic dielectric functions such as gold or bismuth), due to the required computation time, as reported in Refs. \cite{castellanos2021phdthesis, castellanos2023theory}. Therefore, the methods developed previously only allow for the calculation of AMT in small nanoparticles, up to $a=5$ nm.


It has been demonstrated that there are terms in the Maxwell stress tensor, corresponding to the external fields produced by the electron, that do not contribute to the total AMT \cite{castellanos2023theory}. Therefore, it has been shown that the integral of the Maxwell stress tensor containing only the electron's electromagnetic fields must vanish. The interaction term in the stress tensor, which includes both the electron's electromagnetic fields and the fields scattered by the nanoparticle, contributes the most to momentum transfer, and the term including only the fields scattered by the nanoparticle, although small, does not vanish \cite{castellanos2021phdthesis, castrejon2021phdthesis}.
\section{General objective and specific objectives}


The ${\color{red}\text{general objective}}$ of this work is to develop a new methodology for the theoretical study of angular momentum transfer from a swift electron AMT to a spehrical NP, using a classical electrodynamics approach. 

Among the key ${\color{red}\text{specific objectives}}$, the first one is to derive an analytical solution for the integrals of the Maxwell stress tensor in frequency domain. Previously, these integrals were computed numerically using the cubature method.\footnote{The quadrature method is different from the cubature method. Cubatures are used to numerically solve multidimensional integrals \cite{hahn2007cuba}.} The analytical solution will provide an efficient and systematic study of fast electron AMT to NPs with radii up to 50 nm, thus covering the entire nanoscale. In \hyperref[cap: teoria metodos]{Section} \ref{cap: teoria metodos}, the theoretical framework and methods necessary for deducing the electromagnetic fields generated by a relativistic electron in uniform rectilinear motion are developed. Additionally, the scattered electromagnetic fields by a NP centered at the origin, whose electromagnetic response is characterized by its dielectric function $\epsilon(\omega)$, are discussed. The AMT of the fast electron to the NP is then derived through an integral of the Maxwell stress tensor in frequency space, based on the principle of angular momentum conservation in electrodynamics.

The second ${\color{red}\text{specific objective}}$ is to apply the analytical solution to study a variety of materials, implementing a numerical code that computes the frequency-domain integral using quadrature methods. For further details of the code, the reader may visit Ref. \cite{Briseno2024AMTgithub}. In \hyperref[cap: resultados]{Section} \ref{cap: resultados}, the preliminary results are presented, showcasing the calculated solution to the Maxwell stress tensor integral in frequency space, thereby enabling efficient AMT computation.

Lastly, the final ${\color{red}\text{specific objective}}$ is to extend the AMT calculations across the entire nanoscale for different materials. To date, such calculations have primarily focused on Drude-like materials. This work represents the first time that the AMT has been calculated using dielectric functions of more realistic materials, marking a significant advancement in the field.